\documentclass[11pt]{article}
\usepackage[margin=1in]{geometry}
\usepackage{amsmath,amssymb,amsthm}
\title{Disproof of Conjecture 00000005244}
\author{earthking11}
\date{October 6, 2026}
\newtheorem{proposition}{Proposition}
\begin{document}
\maketitle
\section*{The two filed clauses}
Write $c(g)$ for the number of irreducible components at genus $g$.  Under the
terminology of the file, irreducibility means $c(g)=1$.

The first clause says that the moduli space is always irreducible for genus at
least two:
\[
 \forall g\ge2,\qquad c(g)=1.
\]
The next clause says that genus two is exceptional and decomposes into two
irreducible components:
\[
 c(2)=2.
\]

\begin{proposition}
No component-count function can satisfy both clauses.
\end{proposition}
\begin{proof}
Because $2\ge2$, substituting $g=2$ into the universal first clause gives
$c(2)=1$.  The second clause gives $c(2)=2$.  Thus $1=2$, a contradiction.
\end{proof}

The later assertions about the component dimensions, their intersection, and
the discriminant degree do not affect this contradiction.  This disproof is
about the conjunction as filed; it does not choose which of the two mutually
incompatible intended statements should replace it.  A syntactic repair would
at least have to change ``genus at least $2$'' to a range excluding genus $2$,
or remove the claimed genus-two decomposition.

The accompanying Lean theorem quantifies over every possible component-count
function and proves that the conjunction is impossible.
\end{document}
